\section{Link Layer \& Local Area Networks (LAN)}
\subsection{Shared Medium Intuition}
\begin{itemize}
	\item Nodes broadcast on same channel; only one can transmit at a time.
	\item Simultaneous transmit $\Rightarrow$ collision $\Rightarrow$ data lost.
	\item Sender uses full link bandwidth while transmitting.
\end{itemize}

\subsection{Definitions}
\begin{itemize}
	\item \textbf{Bandwidth:} max data rate of channel ($B$).
	\item \textbf{Throughput:} fraction of bandwidth used to send traffic.
	\item \textbf{Goodput:} fraction of bandwidth used for successful data.
	\item $\eta=\frac{\text{avg effective bandwidth}}{\text{total bandwidth}}=\frac{\text{avg effective time}}{\text{total time}}$, $0\le\eta\le1$.
\end{itemize}

\subsection{Multiple Access Domains}
\begin{itemize}
    \item \textbf{Multiple Access Domain:} set of nodes sharing a common channel.
    \item \textbf{Channel Partitioning:} divide channel into smaller "pieces" (time/freq/code).
    \item \textbf{Random Access:} transmit at will; handle collisions (ALOHA, CSMA).
    \item \textbf{Taking Turns:} nodes take turns to transmit (polling/token passing).
\end{itemize}
\subsection{ALOHA protocol}
\begin{itemize}
	\item Random access on shared medium; collisions if overlap. Frame time $T$ (size $L$, rate $B$): $T=L/B$.
	\item \textbf{Binomial model:} $N$ users, each transmits in a slot w.p. $p$.
	\item \textbf{Poisson model:} aggregate attempts are Poisson with rate $G$ (attempts/slot).
	\item \textbf{Slotted ALOHA (align to slots):}
	\begin{itemize}
		\item Success if exactly one tx in a slot.
		\item Binomial: $P(\text{succ})=Np(1-p)^{N-1}$, max at $p=1/N$ gives $S_{\max}\approx 1/e$.
		\item Poisson: $S=Ge^{-G}$, $S_{\max}=1/e$ at $G=1$.
	\end{itemize}
	\item \textbf{Pure ALOHA (no slots):}
	\begin{itemize}
		\item Vulnerable period $2T$.
		\item Binomial (approx): $P(\text{succ})\approx Np(1-p)^{2(N-1)}$, max at $p\approx 1/(2N)$ gives $S_{\max}\approx 1/(2e)$.
		\item Poisson: $S=Ge^{-2G}$, $S_{\max}=1/(2e)$ at $G=1/2$.
	\end{itemize}
\end{itemize}

\subsection{CSMA (Carrier Sense Multiple Access)}
\begin{itemize}
    \item Listen before tx; if channel idle, transmit entire frame; else wait.
    \item \textbf{Propagation delay} $\tau$: time for signal to propagate between 2 farthest nodes.
    \item \textbf{CSMA/CD (Collision Detection):} abort tx on collision detection; reduces wasted time.
    \item \textbf{CSMA/CD efficiency:} $\eta=\frac{1}{1+5\frac{d_{\text{prop}}}{d_{\text{trans}}}}$.
\end{itemize}

\subsubsection{CSMA/CD (Collision Detection)}
\begin{itemize}
	\item Detect collisions while transmitting; send jam; then backoff and retry.
	\item \textbf{Binary exponential backoff:} after $m$-th collision pick $k\in\{0,\ldots,2^m-1\}$, wait $k$ slots.
\end{itemize}
\subsubsection{CSMA/CA (Collision Avoidance)}
\begin{itemize}
	\item Used in Wi-Fi: cannot detect collisions reliably, so avoid them.
	\item If channel idle, wait DIFS then transmit; if busy, pick random backoff slots.
	\item Receiver replies with ACK after SIFS; missing ACK implies collision.
\end{itemize}
\subsubsection{CSMA persistence}
\begin{itemize}
	\item \textbf{1-persistent:} transmit immediately when channel becomes idle (high collision under load).
	\item \textbf{0-persistent:} wait random backoff after busy (lower collision, higher delay).
	\item \textbf{$p$-persistent:} when idle, transmit with prob. $p$ each slot (tradeoff).
\end{itemize}

\subsection{Single (Logical) Link in practice}
\subsubsection{MAC Addresses}
\begin{itemize}
	\item Link-layer address unique per NIC (Network Interface Card); assigned by IEEE OUI.
	\item Used to deliver frames on same LAN; flat (non-hierarchical), portable across LANs.
	\item Usually burned-in, but can be software-set.
\end{itemize}
\subsubsection{Ethernet}
\begin{itemize}
	\item Dominant wired LAN; many speed generations; cheap NICs.
	\item \textbf{CSMA/CD (classic shared medium):}
	\begin{enumerate}
		\item Sense channel; if idle transmit, else wait.
		\item If collision detected while tx, abort and send jam.
		\item \textbf{Binary exponential backoff:} after $m$-th collision pick $k\in\{0,\ldots,2^m-1\}$, wait $k\cdot512$ bit-times, retry.
	\end{enumerate}
\end{itemize}

\subsection{Interconnecting broadcast domains}
\begin{itemize}
	\item \textbf{Switch:}
    \begin{itemize}
        \item link-layer device that breaks broadcast domains; stores/forwards frames using MAC table.
        \item Transparent to hosts; self-learning (no manual config), entries age out (TTL).
        \item Each link is its own collision domain; enables simultaneous transmissions.
        \item Forwarding: if dest known on incoming port $\Rightarrow$ drop; else forward to port; if unknown $\Rightarrow$ flood.
		\item \textbf{Switch table:} (MAC, port, TTL). Learn from source MAC+ingress port on each frame.
    \end{itemize}
    \item \textbf{Loops:} cause broadcast storms, multiple frame copies, MAC table instability.
	\item \textbf{Spanning Tree Protocol (STP):} disable links to eliminate loops; elect root, compute shortest-path tree.
	\item \textbf{HELLO/BPDU fields:} SID (switch ID), RID (root ID), DTR (distance to root).
	\item Root selection: prefer lower RID; then lower DTR; then lower SID (tie-break).
\end{itemize}

\subsection{Errors Detection}
\begin{itemize}
	\item Goal: detect (and sometimes correct) bit errors introduced by noise on the link.
	\item \textbf{Repetition code (3-repetition):} send each bit 3 times; majority vote corrects 1 error, detects up to 2.
	\item \textbf{1D parity:} add one parity bit so total 1s is even/odd; detects any odd number of bit errors, cannot locate/correct.
	\item \textbf{2D parity:} arrange bits in a matrix with row+column parity; detects many multi-bit errors and can correct a single-bit error (intersection).
	\item \textbf{CRC:} \todo \textbf{Consider removing this if not in Exams} treat bits as polynomial over $GF(2)$; append $r$-bit remainder so frame is divisible by generator $G(x)$.
	\item CRC detects all single-bit errors, all burst errors of length $\le r$, and most longer bursts (prob. of miss $\approx 2^{-r}$).
\end{itemize}

\section{"Layer 2.5" - ARP}
ARP is generally considered to be part of Layer 2 (Data Link Layer), but it works directly to support Layer 3 (Network Layer).
\begin{itemize}
    \item Why it's Layer 2: ARP messages are encapsulated directly into Ethernet frames. They do not have an IP header (Layer 3) and they cannot be routed outside of your local network (LAN).
    \item Why it feels like Layer 3: Its entire purpose is to handle Layer 3 information (IP addresses). It is the mechanism that allows Layer 3 to actually function on physical hardware.
\end{itemize}
\subsection{Address Resolution Protocol (ARP)}
\begin{itemize}
	\item Translates between MAC addresses and IP addresses on a LAN.
	\item Needed because Layer 3 provides the IP, but frames on Layer 2 need a destination MAC.
	\item To send outside the LAN, the packet uses the IP of the destination but the MAC of the next-hop router; routers then forward using their own L2 on each hop.
	\item Hosts use DHCP and routing to know whether the destination IP is on the same LAN or remote.
	\item Each node maintains an \textbf{ARP cache} (IP $\rightarrow$ MAC table).
	\item If the mapping is missing, send an \textbf{ARP request} (broadcast in the local broadcast domain).
	\item The target replies with an \textbf{ARP response} containing its MAC, which is stored in the cache.
	\item ARP messages stay within a single broadcast domain (not routed).
\end{itemize}